\chapter{Production failure modes in UI runtimes}
\label{ch:failure-modes}

Most UI runtime bugs are not "wrong output" bugs.
They are \emph{systems} bugs: feedback loops, retained graphs, scheduling pathologies, and ownership confusion.

This chapter collects failure modes that show up across runtimes (VDOM, fine-grained reactive, compiler-assisted).
The goal is not fear.
The goal is to give you a checklist you can apply while designing, profiling, and debugging.

\section{Infinite update loops}

\textbf{Symptom:} the runtime never settles; CPU spikes; the UI starves input.

\textbf{Root causes:}
\begin{itemize}
  \item an effect writes to a signal it depends on without a guard
  \item patching triggers synchronous callbacks that enqueue more work
  \item scheduling policy promotes work aggressively without yielding
\end{itemize}

\textbf{Mitigations:}
\begin{itemize}
  \item enforce a "no self-write" discipline (or require explicit \texttt{untrack})
  \item guard effects with equality checks
  \item use budgets + lanes so the runtime yields under load
\end{itemize}

\section{Zombie effects (effects that outlive their owners)}

\textbf{Symptom:} removed components still react to updates; unexpected DOM writes; memory grows.

\textbf{Root causes:}
\begin{itemize}
  \item effect cleanup not called on unmount
  \item dependency edges not detached on disposal
  \item shared global scheduler keeps references to disposed computations
\end{itemize}

\textbf{Mitigations:}
\begin{itemize}
  \item make disposal idempotent and test it
  \item ensure the reactive graph detaches edges on cleanup
  \item treat "ownership" as a first-class runtime concept
\end{itemize}

\section{Unbounded microtasks / event loop starvation}

\textbf{Symptom:} UI becomes unresponsive even though each unit of work is small.

\textbf{Root causes:}
\begin{itemize}
  \item scheduling uses microtasks for flush and re-enqueues itself repeatedly
  \item promise chains prevent macrotask yielding
  \item a flush loop drains all work without a budget
\end{itemize}

\textbf{Mitigations:}
\begin{itemize}
  \item enforce time budgets and chunk work across frames
  \item promote fairness across lanes (input, sync, transition, idle)
  \item allow an explicit yield boundary
\end{itemize}

\section{Memory retention}

\textbf{Symptom:} steady memory growth after navigation or repeated mounts.

\textbf{Typical retention paths:}
\begin{itemize}
  \item event listeners not removed
  \item DOM node references stored in VNodes that stay reachable
  \item reactive graph edges keep computations alive
  \item caches (key maps / memo tables) never evict
\end{itemize}

\textbf{Mitigations:}
\begin{itemize}
  \item teardown tests that assert listeners are removed and effects disposed
  \item strict "owner owns children" rules for cleanup
  \item instrumentation counters (allocations, live nodes, live effects)
\end{itemize}

\section{Incorrect DOM ownership}

\textbf{Symptom:} double-removal bugs, detached nodes that still receive writes, hydration mismatches.

\textbf{Root causes:}
\begin{itemize}
  \item two subsystems believe they own the same DOM subtree
  \item compiled blocks are embedded without clear lifecycle boundaries
  \item hydration attaches to markup that does not match template expectations
\end{itemize}

\textbf{Mitigations:}
\begin{itemize}
  \item make ownership boundaries explicit (mount returns a disposable handle)
  \item ensure destroy is idempotent and verified by tests
  \item treat hydration mismatch as a controlled bailout, not undefined behavior
\end{itemize}

\section{A practical checklist}

When you design or review a runtime subsystem, ask:

\begin{itemize}
  \item Can work become self-sustaining (infinite loop)?
  \item When something is removed, does \emph{all} associated work stop?
  \item Under load, does the runtime yield fairly?
  \item Can the system prove cleanup (listeners, effects, references)?
  \item Is DOM ownership unambiguous at every boundary (VDOM \(\leftrightarrow\) HRBR, hydration, fallback)?
\end{itemize}
